\documentclass[10pt,a4paper]{amsart}
\usepackage[margin=19mm]{geometry}
\usepackage{amsmath,amssymb,mathtools}
\usepackage[scr=rsfs,cal=euler]{mathalfa}
\usepackage{xfrac}
\usepackage[unicode,hidelinks,bookmarks=false]{hyperref}
\newcommand{\ie}{i.e.\ }
\newcommand{\cf}{cf.\ }
\newcommand{\resp}{respectively\ }
\newcommand{\Subsection}{\subsection}
\providecommand{\nobreakdash}{\nobreak\hbox{-}}
\setlength{\emergencystretch}{2em}
\multlinegap0pt
\usepackage{amssymb}
\usepackage{mathtools}
\usepackage{bm}
\usepackage{enumerate}
\usepackage{float}
\usepackage{xcolor}
\usepackage[normalem]{ulem}
\newtheorem{prop}{Proposition}
\newcommand{\GG}{{\mathbb G}}
\newcommand{\PP}{{\mathbb P}}
\DeclareMathOperator{\chara}{char}
\title{Scrollar invariants of singular curves on toric surfaces}
\author{Karl Christ, Xiang He,, Ilya Tyomkin}
\date{Source: arXiv:2608.14080v1 (CC BY 4.0)}
\begin{document}
\maketitle

\noindent\emph{Source and licence.} Scrollar invariants of singular curves on toric surfaces. Version arXiv:2608.14080v1 (CC BY 4.0), distributed by arXiv under the Creative Commons Attribution 4.0 International licence, http://creativecommons.org/licenses/by/4.0/, as recorded in the arXiv licence metadata for this exact version and shown on the abstract page https://arxiv.org/abs/2608.14080v1. Full licence notice: \texttt{licenses/arxiv-2608.14080-CC-BY-4.0.txt}.\par\smallskip

\noindent\textit{Editorial scope.} Counted unit: the article's prop prop:primratcovwithprofile (source0.tex lines 393--395). The article's complete original proof of it is delivered with it (source0.tex lines 397--413, one \texttt{proof} environment of 17 lines). No mathematics is written, repaired or re-derived: the statement and the proof are the article's own lines, in the article's own notation, and no retained line is substituted or re-flowed.  No further source-local context is needed: every cross-reference of every retained line resolves inside the two delivered spans.  Nothing was omitted from any retained span, so the package carries no omission record.  No retained line contains a citation, so the wrapper carries no bibliography and no bibliography entry.  No retained line contains a figure environment, an image-inclusion command or drawing code, so no image is delivered and the package declares no asset dependency.  Editorial formatting only, in addition: an AMS \texttt{amsart} preamble that restates the article's own declarations and macros used by the retained lines, namely the environments \texttt{prop} on separate counters, together with the standard packages those lines need.  The retained environments print in this document as Proposition 1 in order of appearance, whereas the article prints the same items with the numbering of its own full document.  The complete native source of the article as distributed in the arXiv e-print for this exact version is bundled unchanged as \texttt{sources/207632/source0.tex}.
\begin{prop}\label{prop:primratcovwithprofile}
    Let $\Pi_0, \Pi_\infty$ be two partitions of a positive integer $d$ such that \mbox{$\gcd(\Pi_0\cup \Pi_\infty)=1$}. Then, the general endomorphisms of $(\PP^1;0,1,\infty)$ of degree $d$, whose local degrees at the preimages of $0$ and $\infty$ are prescribed by $\Pi_0$ and $\Pi_\infty$, respectively, are primitive. 
\end{prop}

\begin{proof}
    The proof is a straightforward dimension count. The space of endomorphisms satisfying the assumptions of the proposition is irreducible of dimension $|\Pi_0|+|\Pi_\infty|-2$, where $|\Pi|$ denotes the number of elements in the partition $\Pi$. 
    
    Assume that a general endomorphism $f$ is not primitive. Then there exists a factorization $f=f_2\circ f_1$, where $f_i$ are endomorphisms of $(\PP^1;0,1,\infty)$ of degrees $1<d_i<d$. Furthermore, by a suitable choice of coordinates on the intermediate $\PP^1$, we may assume that $|f_1^{-1}(0)|\le |f_1^{-1}(p)|$ for any $p\in f_2^{-1}(0)$, and similarly for $f_1^{-1}(\infty)$. In particular, 
    $$|f^{-1}(0)|\ge |f_1^{-1}(0)|\cdot |f_2^{-1}(0)|\quad \text{and}\quad |f^{-1}(\infty)|\ge |f_1^{-1}(\infty)|\cdot |f_2^{-1}(\infty)|.$$ 
    
    On the other hand, since $f$ is general, the dimension of its parameter space cannot exceed the dimension of the locus of such compositions $f_2 \circ f_1$. Bounding the dimensions of the parameter spaces for $f_1$ and $f_2$ strictly by their profiles over $0$ and $\infty$, we conclude that 
    \begin{equation}\label{eq:dimineq}
        |f^{-1}(0)|+|f^{-1}(\infty)|-2\le |f_1^{-1}(0)|+|f_1^{-1}(\infty)|-2 +|f_2^{-1}(0)|+|f_2^{-1}(\infty)|-2.
    \end{equation}
    It follows that all inequalities above are equalities and
    $$(|f_1^{-1}(0)|-1)(|f_2^{-1}(0)|-1)=(|f_1^{-1}(\infty)|-1)(|f_2^{-1}(\infty)|-1)=0.$$

    If $|f_2^{-1}(0)|=|f_2^{-1}(\infty)|=1$ then $\gcd(\Pi_0\cup \Pi_\infty)$ is divisible by $\deg(f_2)$, which is a contradiction. Thus, without loss of generality, $|f_2^{-1}(0)|>1=|f_1^{-1}(0)|$ and each of the points in $f_2^{-1}(0)$ has a unique preimage under $f_1$. Plainly, there exists at most one such endomorphism $f_1$, and therefore
    $$|f^{-1}(0)|+|f^{-1}(\infty)|-2=|f_2^{-1}(0)|+|f_2^{-1}(\infty)|-2.$$
    Since \eqref{eq:dimineq} is an equality, it follows now that $|f_1^{-1}(\infty)|=1$, and therefore, $f_1(t)=t^{\deg(f_1)}$. If $\chara(K)\mid\deg(f_1)$ then $\chara(K)\mid\gcd(\Pi_0\cup \Pi_\infty)$, which is a contradiction. Otherwise, $f_1$ is unramified over $\GG_m$, which contradicts the fact that each of the points in $f_2^{-1}(0)$ has a unique preimage under $f_1$.
\end{proof}

\end{document}
